\chapter{Langage, fonctions et calcul : les bases}
\section*{Exercice 1.1 -- Nier une affirmation}
[Base] Nier : pour tout réel $x$, il existe un entier $n$ tel que $n>x$. L'affirmation initiale est-elle vraie ?

\section*{Exercice 1.2 -- Domaine avant calcul}
[Base] Déterminer le domaine de $f(x)=\ln(x-1)/\sqrt{4-x}$ puis résoudre $f(x)=0$.

\section*{Exercice 1.3 -- Une équivalence à corriger}
[Raisonnement] Résoudre $\sqrt{x+6}=x$. Expliquer pourquoi une élévation au carré seule ne suffit pas.

\section*{Exercice 1.4 -- Composer}
[Base] Pour $f(x)=x^2$ et $g(x)=x+1$, calculer $g\circ f$ et $f\circ g$. Sont-elles égales ?

\section*{Exercice 1.5 -- Bijection et restriction}
[Raisonnement] Donner un domaine et un ensemble d'arrivée pour rendre $x\mapsto x^2$ bijective, puis écrire sa réciproque.

\section*{Exercice 1.6 -- Égalité d’ensembles}
[Preuve] Pour $f:E\to F$ et $A,B\subset F$, montrer $f^{-1}(A\cup B)=f^{-1}(A)\cup f^{-1}(B)$. Ici $f^{-1}(A)$ désigne une image réciproque, même sans bijection.
